% © 2026 Ing. David Jaroš — CC BY-NC-ND 4.0
%
% This work is licensed under a Creative Commons Attribution-NonCommercial-NoDerivatives
% 4.0 International License (CC BY-NC-ND 4.0).
%
% License History: Earlier drafts (up to v0.3) were released under CC BY 4.0.
% From v0.4 onward, all material is released under CC BY-NC-ND 4.0 to protect
% the integrity of the theoretical work during ongoing academic development.
%
% See LICENSE.md for full license text.

\documentclass[12pt,a4paper]{article}
\usepackage[a4paper,margin=2.5cm]{geometry}
\usepackage{amsmath,amssymb,amsthm}
\usepackage{booktabs}
\usepackage{hyperref}
\usepackage{tcolorbox}
\tcbuselibrary{skins,breakable}

\newtheorem{theorem}{Theorem}[section]
\newtheorem{proposition}[theorem]{Proposition}
\theoremstyle{remark}
\newtheorem{remark}[theorem]{Remark}

\newcommand{\CC}{\mathbb{C}}
\newcommand{\HH}{\mathbb{H}}
\newcommand{\ZZ}{\mathbb{Z}}

\title{\textbf{Gap C2 Step 1: SU(3) Colour-Charge Lattice Attempt}\\
\large Colour triplet structure and the target constraint $Y\in \frac{1}{3}\ZZ$}
\author{Ing.\ David Jaroš}
\date{2026-05-12}

\begin{document}
\maketitle

\section*{Scope}
This note isolates \textbf{Step 1} of Gap C2:
\[
\mathrm{SU}(3)_c\ \text{fundamental representation}
\ \Longrightarrow\
Y\in \tfrac{1}{3}\ZZ.
\]
Only algebraic and representation-theoretic inputs from active canonical files are used.

\section{Inputs used}
\begin{itemize}
  \item \textbf{[L0]} $\mathrm{SU}(3)_c$ acts on quarks in the fundamental representation $\mathbf{3}$ with basis $(e_1,e_2,e_3)$ and generators $\lambda_a/2$; see \texttt{canonical/interactions/sm\_gauge.tex §G.B}.
  \item \textbf{[L0]} $\ZZ_2\times \ZZ_2\times \ZZ_2$ involution structure yielding three colour sectors; see \texttt{canonical/su3\_derivation/su3\_from\_involutions.tex}.
  \item \textbf{[L1]} Gell-Mann--Nishijima relation $Q=T_3+Y/2$; see \texttt{canonical/interactions/qed.tex}.
\end{itemize}

\section{Triplet-sum integrality proposition}
\begin{proposition}[Colour-triplet sum rule under singlet integrality hypothesis]
\label{prop:triplet_sum}
Let $\Phi_{\mathrm{color}}=(\Phi_r,\Phi_g,\Phi_b)$ be a colour-triplet field with common hypercharge eigenvalue $y$.
If the colour-antisymmetric singlet built from three triplet components has integer electromagnetic charge, then
\[
3y\in\ZZ,
\qquad\text{hence}\qquad
y\in\tfrac{1}{3}\ZZ.
\]
\end{proposition}

\begin{proof}
For a colour triplet with colour-independent hypercharge eigenvalue $y$, the hypercharge contribution to the three-colour singlet sum is $y+y+y=3y$.
If the colour-singlet electromagnetic charge is required to be integer and all non-hypercharge terms in $Q=T_3+Y/2$ contribute integer/half-integer values that cancel in the singlet channel, then the remaining hypercharge contribution must satisfy $3y\in\ZZ$.
Therefore $y\in \frac{1}{3}\ZZ$.
\end{proof}

\begin{remark}
Proposition~\ref{prop:triplet_sum} is a formal implication once the singlet-integrality hypothesis is available from UBT as a theorem.
The unresolved point is deriving that hypothesis purely from $\CC\otimes\HH$ and $S[\Theta]$ without importing empirical hadron charge assignments.
\end{remark}

\section{Target theorem status}
\begin{theorem}[SU(3) colour lattice forcing $Y\in\frac{1}{3}\ZZ$]
\label{thm:colour_lattice}
The requirement that the $\mathrm{SU}(3)_c$ colour gauge group in $\CC\otimes\HH$ acts consistently on the quark fundamental representation $\mathbf{3}$ forces hypercharge eigenvalues of colour-triplet fields to lie in $\frac{1}{3}\ZZ$.
\end{theorem}

\begin{proof}[Honest status]
The proof is \textbf{OPEN}: current files establish the triplet structure and Proposition~\ref{prop:triplet_sum}, but the crucial singlet-integrality premise is not yet closed as a standalone theorem from $\CC\otimes\HH$ alone.
Therefore Theorem~\ref{thm:colour_lattice} is currently conditional, not closed at unconditional [L1].
\end{proof}

\begin{tcolorbox}[title={Gap C2 Step 1 — Final verdict}]
Krok 1: SU(3) $\rightarrow Y\in\frac{1}{3}\ZZ$: \textbf{[OPEN — blocker: singlet-integrality theorem from $\CC\otimes\HH$ is missing]}\\
Krok 2: SU(2)$_L \rightarrow Y\in\frac{1}{2}\ZZ$: \textbf{[PROVED L1]} (Gap C1)\\
Krok 3: gcd$(1/3,1/2)=1/6 \rightarrow Y_Q=1/6$: \textbf{[PROVED L0]} (modular arithmetic)\\
Gap C2 celkem: \textbf{[CONDITIONAL]}.
\end{tcolorbox}

\end{document}
